\documentclass{article}
\usepackage[paperwidth=960pt,paperheight=540pt,margin=45pt]{geometry}
\usepackage{fontspec,xeCJK,tikz,xcolor,hyperref}
\setmainfont{TeX Gyre Heros}
\setCJKmainfont{Noto Sans CJK SC}
\definecolor{bg}{HTML}{102B30}
\definecolor{fg}{HTML}{E9F2E9}
\definecolor{acc}{HTML}{72CBC0}
\definecolor{muted}{HTML}{A7C3C5}
\hypersetup{pdftitle={AURORA Alpha Harness},pdfauthor={AURORA Research Team},pdfsubject={可审计的自演化量化研究智能体}}
\pagestyle{empty}
\begin{document}
\begin{tikzpicture}[remember picture,overlay]
\fill[bg] (current page.south west) rectangle (current page.north east);
\draw[acc!30,line width=.5pt] ([xshift=25pt,yshift=25pt]current page.south west) rectangle ([xshift=-25pt,yshift=-25pt]current page.north east);
\foreach \i in {0,...,3}{\draw[acc!15,line width=1pt] ([xshift=740pt,yshift=-80pt-\i*30pt]current page.north west)--++(100pt,-25pt)--++(-100pt,-25pt)--++(-100pt,25pt)--cycle;}
\end{tikzpicture}
\color{fg}
{\large\color{muted}机制驱动 · 实验有界 · 证据可追溯}\par
\vspace{34pt}
{\fontsize{64}{70}\selectfont\bfseries AURORA}\par
\vspace{8pt}
{\fontsize{26}{32}\selectfont ALPHA HARNESS}\par
\vspace{24pt}
{\fontsize{23}{30}\selectfont 可审计的自演化量化研究智能体}\par
\vspace{12pt}
{\large\color{muted}将机制假设、训练线索和研究证据连接为可复核的工作流。}\par
\vspace{32pt}
\begin{minipage}[t]{.29\textwidth}
{\Large\color{acc}01\quad 提出与冻结}\par\medskip
机制合同与受限表达式\par
测量前记录研究对象
\end{minipage}\hfill
\begin{minipage}[t]{.29\textwidth}
{\Large\color{acc}02\quad 检验与演化}\par\medskip
训练线索形成新的假设\par
子代重新接受独立检验
\end{minipage}\hfill
\begin{minipage}[t]{.29\textwidth}
{\Large\color{acc}03\quad 记录与复核}\par\medskip
保存支持、冲突与未知\par
探索不等于正式确认
\end{minipage}
\vfill
{\color{muted}2026-09-09 停止快照：15 个完成候选 · 6 个已评估子代 · 独立确认 0}\par
\medskip
{\small\color{muted}研究原型与探索结果；不构成已验证的交易收益或统计保证。}
\end{document}
